\documentclass[UTF8,a4paper,11pt]{ctexart}
\usepackage[a4paper,margin=2.2cm]{geometry}
\usepackage{amsmath,amssymb,booktabs,enumitem,hyperref}

\title{LaWAM 第一阶段 LAM 损失设计}
\author{}
\date{}

\begin{document}
\maketitle

\section{目标与总损失}

第一阶段训练 Latent Action Model（LAM）：模型从两帧视觉观测中提取低维 latent action，并用它重建未来帧的 DINO patch feature。

\begin{equation*}
\text{video}_{0:T-1}
\rightarrow \text{DINO features}
\rightarrow \text{LAM encoder}
\rightarrow z
\rightarrow \text{LAM decoder}
\rightarrow \widehat{h}_{t+1}.
\end{equation*}

默认 $T=2$、latent dimension 为 $32$。总损失为

\begin{equation}
\mathcal{L}_{\mathrm{total}}
=\mathcal{L}_{\mathrm{recon}}
+\lambda_{\mathrm{div}}\mathcal{L}_{\mathrm{entropy}}
+\mathcal{L}_{\mathrm{VQ/VAE}}
+\lambda_{\mathrm{aux}}\mathcal{L}_{\mathrm{state}}.
\end{equation}

\section{DINO feature reconstruction loss}

DINO encoder 取倒数第二层 \texttt{hidden\_states[-2]}，删除前 5 个特殊 token，保留 spatial patch tokens。对于 $256\times256$ 图像和 patch size $16$，每帧包含 $16\times16=256$ 个 patch token。

当前配置启用

\begin{verbatim}
norm_latents: true
norm_latents_type: "ln"
\end{verbatim}

因此每个 DINO patch token 都经过无 affine 参数的 LayerNorm。decoder 预测下一帧 feature $\widehat{h}_{t+1}$，目标为真实下一帧 feature $h_{t+1}$。

代码支持以下重建损失：

\begin{align*}
\mathcal{L}_{\mathrm{L1}}
&=\operatorname{mean}|\widehat{h}_{t+1}-h_{t+1}|,\\
\mathcal{L}_{\mathrm{L2}}
&=\operatorname{mean}(\widehat{h}_{t+1}-h_{t+1})^2,\\
\mathcal{L}_{\mathrm{SmoothL1}}
&=\operatorname{SmoothL1}(\widehat{h}_{t+1},h_{t+1};\beta=0.1),\\
\mathcal{L}_{\mathrm{cos}}
&=\mathcal{L}_{\mathrm{SmoothL1}}
+1-\operatorname{mean}[\cos(\widehat{h}_{t+1},h_{t+1})].
\end{align*}

当前配置如下：

\begin{center}
\begin{tabular}{lll}
\toprule
配置 & \texttt{loss\_type} & 实际损失 \\
\midrule
\texttt{dino\_large\_vae.yaml} & \texttt{smooth\_l1} & Smooth L1，$\beta=0.1$ \\
\texttt{dino\_base\_vae.yaml} & \texttt{l2} & MSE \\
\bottomrule
\end{tabular}
\end{center}

代码额外记录 cosine similarity 和 L1 error，但它们默认只是监控指标，不参与反向传播。

\section{VAE bottleneck 与 KL loss}

当前配置使用

\begin{verbatim}
vq_type: "vae"
vq_kwargs:
  layer_norm: true
\end{verbatim}

LAM encoder 输出 $n$ 后，VAE 计算

\begin{equation*}
\mu=W_\mu(\operatorname{LN}(n)),\qquad
\log\sigma^2=W_{\log\sigma^2}(\operatorname{LN}(n)),
\end{equation*}

并在训练时采样

\begin{equation*}
z=\mu+\epsilon\odot\sigma,
\qquad \epsilon\sim\mathcal{N}(0,I).
\end{equation*}

KL loss 为

\begin{equation}
\mathcal{L}_{\mathrm{KL}}
=\frac{1}{2}\operatorname{mean}_B
\sum_i\left(\mu_i^2+\exp(\log\sigma_i^2)-1-\log\sigma_i^2\right).
\end{equation}

实际加入总损失的是

\begin{equation}
\mathcal{L}_{\mathrm{VAE}}=\beta\mathcal{L}_{\mathrm{KL}},
\qquad \beta=5\times10^{-5}\quad\text{（代码默认值）}.
\end{equation}

因此 KL 正则通常远小于 reconstruction loss。VAE 模式没有离散 codebook assignment；代码中的 \texttt{vq\_loss} 在这里实际表示加权后的 KL loss。

\section{离散 VQ loss（当前配置不生效）}

只有将 \texttt{vq\_type} 设为 \texttt{vq}、\texttt{nsvq} 或 \texttt{ema} 时，才会使用离散 VQ。令 encoder 输出为 $z_e$，最近的 codebook vector 为 $e_k$：

\begin{align*}
\mathcal{L}_{\mathrm{codebook}}
&=\|e_k-\operatorname{sg}(z_e)\|_2^2,\\
\mathcal{L}_{\mathrm{commit}}
&=\|\operatorname{sg}(e_k)-z_e\|_2^2.
\end{align*}

普通 VQ 的组合形式为

\begin{equation*}
\mathcal{L}_{\mathrm{VQ}}
=\mathcal{L}_{\mathrm{codebook}}
+\beta\mathcal{L}_{\mathrm{commit}}
+\lambda_{\mathrm{orth}}\mathcal{L}_{\mathrm{orth}}.
\end{equation*}

量化结果使用 straight-through estimator，使 reconstruction gradient 能回传给 encoder。当前 \texttt{vq\_type: vae}，所以这些离散 VQ 项均不生效。

\section{Entropy / diversity loss}

离散 VQ 模式下，代码计算单样本 assignment entropy 和整个 batch 的 codebook entropy：

\begin{equation*}
\mathcal{L}_{\mathrm{entropy}}
=\lambda_{\mathrm{sample}}H_{\mathrm{sample}}
-\lambda_{\mathrm{codebook}}H_{\mathrm{codebook}}.
\end{equation*}

其作用是降低单样本 assignment 的不确定性，同时避免整个数据集只使用少量 code。

当前配置为

\begin{verbatim}
vq_type: "vae"
lambda_diversity: 0.0
\end{verbatim}

所以该项恒为零，不影响当前训练。

\section{状态变化辅助损失}

数据中的状态变化定义为

\begin{equation*}
\Delta s=s_{\mathrm{end}}-s_{\mathrm{start}}.
\end{equation*}

StatePredictor 根据 latent action $z$、初始状态 $s_0$ 和 embodiment id 预测 $\widehat{\Delta s}$。输出经过 $\tanh$，范围为 $[-1,1]$。

状态损失支持 L1 和 L2：

\begin{align*}
\mathcal{L}_{\mathrm{state,L1}}
&=\frac{\sum m|\widehat{\Delta s}-\Delta s|}{\sum m},\\
\mathcal{L}_{\mathrm{state,L2}}
&=\frac{\sum m(\widehat{\Delta s}-\Delta s)^2}{\sum m},
\end{align*}

其中 $m$ 是 state mask，用来忽略 padding 和无效状态维度。

large 配置为

\begin{verbatim}
state_loss_type: "l2"
lambda_aux: 1.0
\end{verbatim}

状态损失只对 \texttt{embodiment\_id != 0} 的 robot 样本计算；human 样本跳过该项。

\section{当前 large 配置的实际目标}

\texttt{dino\_large\_vae.yaml} 的关键配置为

\begin{verbatim}
loss_type: "smooth_l1"
state_loss_type: "l2"
lambda_aux: 1.0
lambda_diversity: 0.0
vq_type: "vae"
vq_kwargs:
  layer_norm: true
\end{verbatim}

所以实际损失为

\begin{equation}
\boxed{
\mathcal{L}_{\mathrm{total}}
=\mathcal{L}_{\mathrm{SmoothL1}}(\widehat{h}_{t+1},h_{t+1})
+\mathcal{L}_{\mathrm{state,L2}}
+5\times10^{-5}\mathcal{L}_{\mathrm{KL}}
}
\end{equation}

其中状态项只对 robot 样本的有效状态维度生效。

\begin{center}
\begin{tabular}{lll}
\toprule
损失项 & 当前是否生效 & 说明 \\
\midrule
DINO reconstruction & 是 & Smooth L1，$\beta=0.1$ \\
VAE KL & 是 & 默认权重 $5\times10^{-5}$ \\
State auxiliary & 有条件 & 仅 robot 样本和有效状态维度 \\
Entropy/diversity & 否 & 权重为零，且 VAE 返回零 \\
离散 VQ loss & 否 & 当前使用 VAE \\
\bottomrule
\end{tabular}
\end{center}

\section{双视图增强}

large 配置还启用了

\begin{verbatim}
image_aug: true
dual_view_aug: true
\end{verbatim}

训练时 encoder 和 decoder target 可以来自同一视频的两套独立增强视图：

\begin{equation*}
\text{view}_1\rightarrow\text{LAM encoder},\qquad
\text{view}_2\rightarrow\text{decoder input and target}.
\end{equation*}

增强包括 RandomResizedCrop 和 ColorJitter。因此模型同时被要求学习一定的视觉增强不变性；这会增加逐 patch feature reconstruction 的难度，也可能提高该损失的可达到下限。

\section{相关实现文件}

\begin{itemize}[nosep]
  \item 总损失与 reconstruction loss：\texttt{latent\_action\_model/core/lam\_lightinng.py}
  \item LAM 数据流与目标 feature：\texttt{latent\_action\_model/core/lam\_model.py}
  \item DINO feature 提取：\texttt{latent\_action\_model/core/vjepa\_encoder.py}
  \item VAE、VQ、entropy loss：\texttt{latent\_action\_model/core/vq.py}
  \item 状态辅助损失：\texttt{latent\_action\_model/core/utils/utils.py}
  \item 当前配置：\texttt{latent\_action\_model/config/dino\_large\_vae.yaml}
\end{itemize}

\end{document}
